\section{Appendix to Section \ref{sec:extensions}}\label{app:extensions}

\subsection{The merger-attempt condition under a clawback}

The clawback introduced in Section \ref{sec:extensions} is modeled as a reduced-form tax on the entrepreneur's merger-related continuation payoff.
It therefore changes the division of gains from acquisition, but not the joint surplus generated by attempting a merger.

Accordingly, the merger-attempt condition remains
\begin{equation}
    q\Delta_j \ge k.
\end{equation}
Equivalently, the attempt indicator is unchanged:
\begin{equation}
    a_j(q)=\mathbf{1}\{q\Delta_j \ge k\}.
\end{equation}
This distinction matters for welfare accounting.
Because the clawback does not affect product-market outcomes or merger feasibility, the local and national social-value objects \(B_j^L(q)\) and \(B_j^N(q)\) are unchanged.
Only private continuation values and therefore private effort, entry, and project choice are modified.
This assumption corresponds to a setting in which the clawback is imposed on the entrepreneur's realized continuation payoff after bargaining, rather than on the gross joint returns from merger attempt.
If the clawback also reduced the joint surplus from merger attempt, then the attempt threshold itself would shift; the present extension abstracts from that channel in order to isolate the composition effect of taxing entrepreneurial buyout returns.
For that reason, the extension should be read as a proof-of-concept policy-design exercise rather than as a full institutional model of clawback implementation.

\subsection{Project choice without the single-crossing restriction under a clawback}

Define the private payoff difference under a clawback:
\begin{equation}
    \widetilde D(q,\rho)\equiv \widetilde R_B(q,\rho)-\widetilde R_I(q,\rho).
\end{equation}
Using \eqref{eq:Rtilde}, this difference is
\begin{equation}\label{eq:Dqpiececlaw}
    \widetilde D(q,\rho)
    =
    \begin{cases}
        \pi_E^D(B)-\pi_E^D(I), &
        q < \bar q_B, \\
        \pi_E^D(B)-\pi_E^D(I)+(1-\rho)\beta(q\Delta_B-k), &
        \bar q_B \le q < \bar q_I, \\
        \pi_E^D(B)-\pi_E^D(I)+(1-\rho)\beta q(\Delta_B-\Delta_I), &
        q \ge \bar q_I,
    \end{cases}
\end{equation}
where
\begin{equation}
    \bar q_j\equiv \frac{k}{\Delta_j}.
\end{equation}

Because \(\pi_E^D(I)>\pi_E^D(B)\), we have
\begin{equation}
    \widetilde D(q,\rho)<0
    \qquad
    \text{for all } q<\bar q_B.
\end{equation}
In the intermediate region \([\bar q_B,\bar q_I)\), the buyout-oriented project is preferred if and only if
\begin{equation}\label{eq:qtilderho}
    q \ge \widetilde q(\rho)
    \equiv
    \frac{
        \pi_E^D(I)-\pi_E^D(B)+(1-\rho)\beta k
    }{
        (1-\rho)\beta\Delta_B
    }.
\end{equation}
Once both projects are merger-feasible, the buyout-oriented project is preferred if and only if
\begin{equation}\label{eq:qhatappendixrho}
    q \ge \hat q(\rho)
    \equiv
    \frac{
        \pi_E^D(I)-\pi_E^D(B)
    }{
        (1-\rho)\beta(\Delta_B-\Delta_I)
    }.
\end{equation}

Thus the clawback shifts upward both the intermediate-region threshold \(\widetilde q(\rho)\) and the full-feasibility threshold \(\hat q(\rho)\).
Under the maintained single-crossing structure used in the main text, attention can be restricted to the latter.

\subsection{Proof of Proposition \ref{prop:clawbackthreshold}}

Suppose both project types are merger-feasible, so that
\begin{equation}
    \widetilde R_j(q,\rho)
    =
    \pi_E^D(j)+(1-\rho)\beta(q\Delta_j-k),
    \qquad j\in\{I,B\}.
\end{equation}
The entrepreneur prefers project \(B\) if and only if
\begin{equation}
    \widetilde R_B(q,\rho)\ge \widetilde R_I(q,\rho).
\end{equation}
Substituting the expressions above yields
\begin{equation}
    \pi_E^D(B)+(1-\rho)\beta(q\Delta_B-k)
    \ge
    \pi_E^D(I)+(1-\rho)\beta(q\Delta_I-k).
\end{equation}
The \(-(1-\rho)\beta k\) terms cancel, so project \(B\) is preferred if and only if
\begin{equation}
    (1-\rho)\beta q(\Delta_B-\Delta_I)
    \ge
    \pi_E^D(I)-\pi_E^D(B).
\end{equation}
Since \(\Delta_B>\Delta_I\), this is equivalent to
\begin{equation}
    q\ge
    \hat q(\rho)
    \equiv
    \frac{
        \pi_E^D(I)-\pi_E^D(B)
    }{
        (1-\rho)\beta(\Delta_B-\Delta_I)
    }.
\end{equation}
This proves \eqref{eq:qhatrho}.

Differentiating \(\hat q(\rho)\) with respect to \(\rho\) gives
\begin{equation}
    \frac{\partial \hat q(\rho)}{\partial \rho}
    =
    \frac{
        \pi_E^D(I)-\pi_E^D(B)
    }{
        (1-\rho)^2\beta(\Delta_B-\Delta_I)
    }
    =
    \frac{\hat q(\rho)}{1-\rho}
    >0.
\end{equation}
This proves \eqref{eq:dqhatdrho}. \qed

\subsection{Derivation of the minimal clawback \eqref{eq:rhobar}}

Suppose \(q\) is such that both project types are merger-feasible.
To make project \(I\) privately optimal, it is sufficient to require
\begin{equation}
    q \le \hat q(\rho).
\end{equation}
Using \eqref{eq:qhatrho}, this becomes
\begin{equation}
    q
    \le
    \frac{
        \pi_E^D(I)-\pi_E^D(B)
    }{
        (1-\rho)\beta(\Delta_B-\Delta_I)
    }.
\end{equation}
Rearranging gives
\begin{equation}
    1-\rho
    \le
    \frac{
        \pi_E^D(I)-\pi_E^D(B)
    }{
        \beta q(\Delta_B-\Delta_I)
    }.
\end{equation}
Hence the smallest clawback that restores project \(I\) is
\begin{equation}
    \bar \rho(q)
    =
    1-
    \frac{
        \pi_E^D(I)-\pi_E^D(B)
    }{
        \beta q(\Delta_B-\Delta_I)
    }.
\end{equation}
Imposing the constraint \(\rho\ge 0\) yields \eqref{eq:rhobar}. \qed

\subsection{Modified welfare and subsidy formulas under a clawback}

Under the clawback extension, private continuation values are given by \(\widetilde R_j(q,\rho)\), and we define
\begin{equation}
    \widetilde A_j(q,\rho)\equiv \frac{\widetilde R_j(q,\rho)^2}{2},
    \qquad
    \widetilde c_j(s,q,\rho)\equiv s+\widetilde A_j(q,\rho).
\end{equation}
Because the clawback does not change realized product-market outcomes or the merger-attempt condition, the social-value objects \(B_j^L(q)\) and \(B_j^N(q)\) remain those defined in Sections \ref{sec:subsidy} and \ref{sec:decentralized}.

For the local planner,
\begin{equation}
    \widetilde W_j^L(s,q,\rho)
    =
    \widetilde c_j(s,q,\rho)
    \left(
        \widetilde R_j(q,\rho)B_j^L(q)
        -
        \widetilde A_j(q,\rho)
        -
        \tau s
    \right)
    -
    \frac{\widetilde c_j(s,q,\rho)^2}{2}.
\end{equation}
Expanding as in Appendix \ref{app:subsidy}, the derivative with respect to \(s\) is
\begin{equation}
    \frac{\partial \widetilde W_j^L(s,q,\rho)}{\partial s}
    =
    \widetilde R_j(q,\rho)B_j^L(q)
    -
    \frac{\tau+2}{2}\widetilde R_j(q,\rho)^2
    -
    (2\tau+1)s.
\end{equation}
Hence the local planner's type-specific optimal subsidy is
\begin{equation}
    \widetilde s_j^{L*}(q,\rho)
    =
    \max\left\{
        0,\;
        \frac{
            \widetilde R_j(q,\rho)B_j^L(q)
            -
            \frac{\tau+2}{2}\widetilde R_j(q,\rho)^2
        }{2\tau+1}
    \right\}.
\end{equation}

Similarly, the national planner's welfare under project \(j\) is
\begin{equation}
    \widetilde W_j^N(s,q,\rho)
    =
    \widetilde c_j(s,q,\rho)
    \left(
        \widetilde R_j(q,\rho)B_j^N(q)
        -
        \widetilde A_j(q,\rho)
        -
        \tau s
    \right)
    -
    \frac{\widetilde c_j(s,q,\rho)^2}{2},
\end{equation}
implying
\begin{equation}
    \widetilde s_j^{N*}(q,\rho)
    =
    \max\left\{
        0,\;
        \frac{
            \widetilde R_j(q,\rho)B_j^N(q)
            -
            \frac{\tau+2}{2}\widetilde R_j(q,\rho)^2
        }{2\tau+1}
    \right\}.
\end{equation}

Whenever both optimal subsidies are interior,
\begin{equation}
    \widetilde s_j^{L*}(q,\rho)-\widetilde s_j^{N*}(q,\rho)
    =
    \frac{
        \widetilde R_j(q,\rho)\bigl(B_j^L(q)-B_j^N(q)\bigr)
    }{2\tau+1}.
\end{equation}
Thus the clawback alters the magnitude of the local--national wedge through private continuation values, while holding fixed the underlying incidence difference between local and national welfare for a given project type.

\subsection{Proof of Proposition \ref{prop:instrumentsubstitution}}

The project-switching condition under a clawback is
\begin{equation}
    q(\theta)=\hat q(\rho).
\end{equation}
Differentiating totally with respect to \(\rho\) yields
\begin{equation}
    q'(\theta)\frac{d\theta}{d\rho}
    =
    \hat q'(\rho).
\end{equation}
Hence
\begin{equation}
    \frac{d\theta}{d\rho}
    =
    \frac{\hat q'(\rho)}{q'(\theta)}.
\end{equation}
Since \(q'(\theta)<0\) and Proposition \ref{prop:clawbackthreshold} implies \(\hat q'(\rho)>0\), it follows that
\begin{equation}
    \frac{d\theta}{d\rho}<0.
\end{equation}
Therefore a higher clawback rate permits a lower antitrust toughness while keeping the entrepreneur on the same project-indifference locus.

The second statement follows immediately from the condition for implementing project \(I\):
\begin{equation}
    q(\theta)\le \hat q(\rho).
\end{equation}
This inequality can be satisfied either by lowering \(q(\theta)\) through stricter antitrust, by raising \(\hat q(\rho)\) through a stronger clawback, or by adjusting both instruments jointly.
This proves the proposition. \qed
